% ============================================================
%  USPSC-Cap3-Metodologia_3.6-Abordagens_de_Classificacao.tex
%  Seção 3.6 — Abordagens de Classificação de Sentimento
% ============================================================

\section{Abordagens de Classificação de Sentimento}
\label{sec:abordagens}

Em consonância com o Objetivo Específico~4 estabelecido na \autoref{sec:objetivos} — avaliar e comparar o desempenho de diferentes abordagens de classificação de sentimento no domínio financeiro em português —, este trabalho avalia três paradigmas metodológicos distintos: a abordagem \textbf{baseada em léxico}, a abordagem \textbf{baseada em modelo de linguagem pré-treinado especializado por domínio} (FinBERT-PT-BR) e a abordagem \textbf{baseada em modelo de linguagem pré-treinado genérico ajustado por \textit{fine-tuning} em dados locais} (BERTimbau). A distinção entre os dois paradigmas neurais é deliberada: enquanto o FinBERT-PT-BR obtém sua especialização por meio de pré-treinamento sobre um corpus financeiro de larga escala em etapa anterior a este trabalho, o BERTimbau parte de um modelo genérico e adquire especialização \textit{a posteriori}, por meio de ajuste fino sobre o próprio corpus de \textit{tweets} financeiros coletados. Essa dicotomia --- especialização prévia de domínio versus especialização \textit{post hoc} de gênero textual e subdomínio --- constitui o eixo central da comparação empírica apresentada no \autoref{Cap:Ava_Exp}. A inclusão de métodos léxicos permite situar ambas as abordagens neurais em relação às ferramentas de referência disponíveis para o português, orientando recomendações práticas para trabalhos futuros \cite{kearney2014,sousa2025}.


\subsection{Abordagem Baseada em Léxico: OpLexicon}
\label{subsec:abordagem_lexica}

Métodos baseados em léxico (\textit{lexicon-based}) classificam o sentimento de um texto a partir da correspondência de seus termos a um dicionário de polaridade previamente construído, no qual cada palavra ou expressão recebe um escore de valência positiva, negativa ou neutra \cite{liu2020survey}. A principal vantagem dessa abordagem reside na transparência e na ausência de necessidade de dados de treinamento rotulados. Sua limitação central está na incapacidade de capturar negação,
ironia e vocabulário específico do domínio não contemplado pelo léxico \cite{loughran2011,kearney2014}.

\subsubsection{Escolha do Recurso Léxico}
\label{subsubsec:escolha_lexico}

O \textbf{OpLexicon v3.0} foi selecionado como léxico principal por duas razões: (i)~abrangência léxica consideravelmente maior (32.191 entradas de \textit{tokens} não lematizados, cobrindo formas coloquiais recorrentes em \textit{tweets}); e (ii)~construção orientada ao português brasileiro, com termos do jargão financeiro local. O \textbf{SentiLex-PT} \cite{carvalho2015sentilex}, por sua vez, é implementado em paralelo como segunda abordagem léxica (Seção~\ref{subsec:abordagem_sentilex}),  permitindo avaliar o impacto da anotação manual e da cobertura de formas flexionadas sobre o desempenho no domínio financeiro.


\subsubsection{Estrutura do Léxico}
\label{subsubsec:estrutura_lexico}

O OpLexicon v3.0 é distribuído como arquivo de texto tabular (\textit{tab-separated values}) sem cabeçalho, com quatro colunas: \texttt{term} (lema), \texttt{type} (classe gramatical), \texttt{polarity} (polaridade inteira: $+1$, $-1$ ou $0$) e \texttt{polarity\_revision} (origem da entrada: \texttt{A} para revisão manual ou \texttt{C} para atribuição automática). A indexação por lemas — e não por formas flexionadas — é a razão pela qual a lematização constitui etapa obrigatória do pré-processamento para esta abordagem, conforme detalhado na \autoref{subsec:fluxo_lexico}.

\subsubsection{Implementação e Protocolo de Pontuação}
\label{subsubsec:protocolo_oplexicon}

A implementação foi realizada no quinto estágio do \textit{pipeline} por meio da classe \texttt{OpLexiconAnalyzer}, definida em \texttt{app/core/processing/lexicon/op\_lexicon.py}. A classe herda de \texttt{BaseSentimentAnalyzer} — a mesma classe abstrata (ABC) que fundamenta o
\texttt{FinBertPTBRAnalyzer} —, garantindo interface uniforme entre as duas abordagens e extensibilidade para classificadores futuros.

O protocolo de pontuação compreende as seguintes etapas:

\begin{enumerate}
	\item \textbf{Seleção dos registros}: o método \texttt{run()} recupera exclusivamente as publicações que já possuem classificação humana em \texttt{tweets\_classification} mas ainda não foram processadas pelo OpLexicon, garantindo que a comparação opere sobre o mesmo
	\textit{gold standard} utilizado pelo FinBERT-PT-BR;
	
	\item \textbf{Pré-processamento}: cada texto é submetido ao método \texttt{preprocess()}, que aplica o fluxo completo descrito na \autoref{subsec:fluxo_lexico}, incluindo lematização via \texttt{spaCy} — etapa obrigatória para maximizar a taxa de correspondência com as entradas do léxico, que indexa lemas e não formas flexionadas;
	
	\item \textbf{Carregamento do léxico}: o método \texttt{load\_model()} lê o arquivo \texttt{lexico\_v3.0.txt} e constrói um dicionário em memória no formato
	\texttt{\{lema: polaridade\}}, com tempo de carga desprezível e sem demanda de GPU;
	
	\item \textbf{Pontuação agregada}: para cada \textit{token} do texto pré-processado, consulta-se o dicionário léxico. \textit{Tokens} ausentes recebem pontuação zero e são excluídos da agregação. A polaridade final é calculada pela média aritmética dos escores dos \textit{tokens} com entrada no léxico, conforme a \autoref{eq:oplexicon_score}:
\end{enumerate}

\begin{equation}
	\label{eq:oplexicon_score}
	\bar{p} \;=\; \frac{1}{|T^{*}|} \sum_{t \,\in\, T^{*}} \mathrm{pol}(t)
\end{equation}

\noindent onde $T^{*} = \{t \in T \mid \mathrm{pol}(t) \neq 0\}$ é o subconjunto de \textit{tokens} com polaridade não-nula e $\mathrm{pol}(t)$ é a polaridade do \textit{token} $t$
no léxico. Caso $T^{*} = \emptyset$, o texto é classificado como Neutro com escore zero.

\begin{enumerate}
	\setcounter{enumi}{4}
	\item \textbf{Limiarização}: a classe final é determinada por um limiar $\theta = 0{,}05$ aplicado à média $\bar{p}$: se $\bar{p} > \theta$, a publicação é classificada como Positivo; se $\bar{p} < -\theta$, como Negativo; caso contrário, como Neutro. O valor reduzido do limiar reflete a alta proporção esperada de \textit{tokens} financeiros sem correspondência no léxico (siglas de ativos, nomes de emissores, jargão de derivativos), que dilui a magnitude da média
	sem necessariamente indicar neutralidade semântica;
	
	\item \textbf{Persistência}: os resultados são armazenados em \texttt{tweets\_classification} com \texttt{classificator = "OpLexicon"} e o respectivo escore, permitindo comparação direta com as anotações humanas e com o FinBERT-PT-BR na etapa de avaliação (\autoref{subsec:metricas_avaliacao}).
\end{enumerate}

% ------------------------------------------------------------------
\subsection{Abordagem Baseada em Léxico: SentiLex-PT}
\label{subsec:abordagem_sentilex}
% ------------------------------------------------------------------

\subsubsection{Fundamentação da Escolha}
\label{subsubsec:justificativa_sentilex}

O \textbf{SentiLex-PT} \cite{carvalho2015sentilex} foi incorporado como segunda abordagem léxica com o objetivo de ampliar a validade comparativa da avaliação. Enquanto o OpLexicon~v3.0 é construído sobre um processo semiautomático voltado ao português brasileiro, o SentiLex-PT decorre de um esforço de anotação manual, o que lhe confere características complementares relevantes para o domínio financeiro:

\begin{enumerate}
	\item \textbf{Anotação manual}: os 7.014 lemas (expandidos para 82.347 formas flexionadas na versão SentiLex-PT02) foram anotados manualmente por especialistas segundo a taxonomia \textit{Manually Annotated} (\texttt{ANOT:MAN}), resultando em um recurso de alta precisão intrínseca \cite{carvalho2015sentilex};

	\item \textbf{Cobertura de formas flexionadas}: a versão SentiLex-PT02, adotada neste trabalho, contém formas nominais, adjetivais e verbais flexionadas, dispensando etapa de lematização e tornando-a mais adequada ao processamento de \textit{tweets} --- textos curtos com vocabulário coloquial e morfologia variada;

	\item \textbf{Validade científica consolidada}: o SentiLex-PT é amplamente utilizado como referência em estudos de análise de sentimento em português \cite{liu2020survey}, permitindo contextualizar os resultados obtidos neste trabalho em relação à literatura.
\end{enumerate}

\subsubsection{Estrutura do Léxico}
\label{subsubsec:estrutura_sentilex}

O SentiLex-PT02 é distribuído como arquivo de texto simples, com uma entrada por linha no formato:

\begin{center}
	\texttt{forma.PoS,POL:\textit{N};ANOT:\textit{TAG}}
\end{center}

\noindent onde \texttt{forma} é a palavra flexionada em letras minúsculas, \texttt{PoS} é a classe gramatical (\texttt{Adj}, \texttt{N}, \texttt{V}, entre outras), \texttt{POL} é a polaridade inteira (\texttt{1}~=~positivo, \texttt{-1}~=~negativo, \texttt{0}~=~neutro) e \texttt{ANOT} é o tipo de anotação. A \autoref{tab:sentilex_exemplos} apresenta exemplos representativos do léxico no domínio financeiro.

\begin{table}[!htbp]
	\caption{Exemplos de entradas do SentiLex-PT02 no domínio financeiro}
	\label{tab:sentilex_exemplos}
	\centering
	\begin{tabular}{lllc}
		\hline
		\textbf{Forma} & \textbf{PoS} & \textbf{ANOT} & \textbf{POL} \\
		\hline
		lucro      & N   & MAN & \phantom{$-$}1 \\
		\hline
		crescimento & N  & MAN & \phantom{$-$}1 \\
		\hline
		valorização & N  & MAN & \phantom{$-$}1 \\
		\hline
		queda      & N   & MAN & $-$1 \\
		\hline
		prejuízo   & N   & MAN & $-$1 \\
		\hline
		crise      & N   & MAN & $-$1 \\
		\hline
		empresa    & N   & MAN & \phantom{$-$}0 \\
		\hline
		mercado    & N   & MAN & \phantom{$-$}0 \\
		\hline
	\end{tabular}
	\fonte{Adaptado de \citeonline{carvalho2015sentilex}.}
\end{table}

\subsubsection{Implementação e Protocolo de Pontuação}
\label{subsubsec:protocolo_sentilex}

A implementação foi realizada no quinto estágio do \textit{pipeline} por meio da classe \texttt{SentiLexAnalyzer}, definida em \texttt{app/core/processing/lexicon/senti\_lex.py}. A classe herda de \texttt{BaseSentimentAnalyzer} (\texttt{app/core/processing/base\_analyzer.py}), garantindo interface uniforme com o \texttt{OpLexiconAnalyzer} e o \texttt{FinBertPTBRAnalyzer}.

O protocolo de pontuação compreende as seguintes etapas:

\begin{enumerate}
	\item \textbf{Seleção dos registros}: o método \texttt{run()} recupera exclusivamente as publicações que já possuem classificação humana em \texttt{tweets\_classification}, utilizadas como \textit{gold standard};

	\item \textbf{Pré-processamento}: cada texto é submetido ao método \texttt{preprocess()}, que aplica o fluxo descrito na \autoref{sec:preprocessamento} --- remoção de URLs, menções, caracteres especiais e conversão para letras minúsculas. Diferentemente do \texttt{OpLexiconAnalyzer}, não é realizada lematização via \texttt{spaCy}, pois o SentiLex-PT02 já contém formas flexionadas;

	\item \textbf{Carregamento do léxico}: o método \texttt{load\_model()} lê o arquivo \texttt{data/SentiLex- PT02.txt} e constrói um dicionário em memória no formato \texttt{\{forma: polaridade\}}, com tempo de carga desprezível e sem demanda de GPU;

	\item \textbf{Pontuação com janela de negação}: para cada \textit{token} do texto pré-processado, consulta-se o dicionário léxico. A fim de capturar negações simples --- padrão frequente em textos financeiros, como \textit{``não houve crescimento''} ---, aplica-se uma janela de negação de três \textit{tokens}: quando um \textit{token} pertencente ao conjunto de marcadores de negação\footnote{Conjunto de marcadores: \textit{não, nunca, jamais, nem, nenhum, nenhuma, tampouco, sequer}.} precede uma palavra com polaridade não-nula dentro desse intervalo, o sinal da polaridade é invertido. A pontuação final é dada por:

	\begin{equation}
		\label{eq:sentilex_score}
		\bar{p} = \sum_{t \in T^{*}} \mathrm{pol\_neg}(t)
	\end{equation}

	\noindent onde $T^{*} = \{t \in T \mid \mathrm{pol}(t) \neq 0\}$ é o subconjunto de \textit{tokens} com polaridade não-nula e $\mathrm{pol\_neg}(t)$ é a polaridade do \textit{token} $t$ após aplicação da janela de negação. Caso $T^{*} = \emptyset$, o texto é classificado como Neutro com escore zero;

	\item \textbf{Limiarização}: a classe final é determinada diretamente pelo sinal da pontuação acumulada $\bar{p}$: se $\bar{p} > 0$, a publicação é classificada como Positivo; se $\bar{p} < 0$, como Negativo; se $\bar{p} = 0$, como Neutro. O escore de confiança reportado é dado por $|\bar{p}| / |T|$, normalizado ao intervalo $[0, 1]$, análogo ao \textit{score} de probabilidade retornado pelo FinBERT-PT-BR;

	\item \textbf{Persistência}: os resultados são armazenados em \texttt{tweets\_classification} com \texttt{classificator = ``SentiLex-PT''} e o respectivo escore, permitindo comparação direta com as demais abordagens conforme descrito na \autoref{subsec:metricas_avaliacao}.
\end{enumerate}


\subsection{Abordagem Baseada em Modelo de Linguagem Pré-treinado: FinBERT-PT-BR}
\label{subsec:finbert}

\subsubsection{Fundamentação da Escolha}
\label{subsubsec:justificativa_finbert}

O \textbf{FinBERT-PT-BR} \cite{santos2023finbert} foi selecionado como ferramenta central da abordagem neural por três razões fundamentais, já antecipadas no \autoref{Cap:Fundamentacao_Teorica}:

\begin{enumerate}
	\item \textbf{Especialização no domínio}: o modelo foi pré-treinado sobre um corpus de 1,4 milhão de textos financeiros em português, incluindo notícias, relatórios e publicações de veículos especializados, o que lhe confere representações contextuais calibradas para o vocabulário técnico do mercado financeiro brasileiro;
	
	\item \textbf{Desempenho superior no estado da arte}: nos experimentos reportados por \citeonline{santos2023finbert}, o FinBERT-PT-BR superou modelos de linguagem genéricos e abordagens léxicas em tarefas de classificação de sentimento financeiro em português;
	
	\item \textbf{Disponibilidade pública}: o modelo está disponível publicamente no repositório \textit{Hugging Face} sob o identificador \texttt{lucas-leme/FinBERT-PT-BR}, garantindo a reprodutibilidade da metodologia proposta.
\end{enumerate}

\subsubsection{Arquitetura do Modelo}
\label{subsubsec:arquitetura_finbert}

O FinBERT-PT-BR é construído sobre a arquitetura \textbf{BERT} (\textit{Bidirectional Encoder Representations from Transformers}) \cite{devlin2019}, que implementa o mecanismo de atenção multi-cabeça proposto por \citeonline{vaswani2017}. A característica central do BERT — a codificação bidirecional do contexto — permite que a
representação de cada \textit{token} incorpore informações tanto dos termos que o precedem quanto dos que o sucedem na sequência, o que é particularmente relevante para a interpretação de negações e modificadores de intensidade frequentes no discurso financeiro (e.g., ``\textit{não} sobe'', ``\textit{queda} acentuada'').

O modelo base utilizado para o \textit{fine-tuning} é o \textbf{BERTimbau} \cite{souza2020}, pré-treinado sobre um corpus de 2,68 bilhões de palavras em português brasileiro. Sobre esse modelo base, \citeonline{santos2023finbert} realizaram o \textit{fine-tuning} para a tarefa de classificação de sentimento financeiro em três classes:
\textbf{Positivo}, \textbf{Negativo} e \textbf{Neutro}.

\subsubsection{Implementação e Protocolo de Inferência}
\label{subsubsec:protocolo_inferencia}

A inferência foi implementada no quinto estágio do \textit{pipeline} (\texttt{app/core/processing /bert/}), por meio da classe \texttt{FinBertPTBRAnalyzer}, definida em \texttt{app/core/processing/ bert/finbert\_ptbr.py}. A classe herda de \texttt{BaseSentimentAnalyzer} (\texttt{app/core/process ing/base\_analyzer.py}), uma classe abstrata (ABC) que estabelece o contrato de interface para qualquer analisador de sentimento no \textit{pipeline}, garantindo extensibilidade para futuros modelos.

O protocolo de inferência compreende as seguintes etapas:

\begin{enumerate}
	\item \textbf{Seleção dos registros}: o método \texttt{run()} recupera via \texttt{app/shared/db/database.py} exclusivamente as publicações que já possuem classificação humana em \texttt{tweets\_classification} mas ainda não foram processadas pelo FinBERT-PT-BR. Esse critério garante que a inferência do modelo opere sobre o mesmo subconjunto utilizado como \textit{gold standard}, permitindo comparação direta entre as duas abordagens;
	
	\item \textbf{Pré-processamento}: cada texto é submetido ao método \texttt{preprocess()} da própria classe, que aplica o fluxo reduzido descrito na \autoref{subsec:fluxo_finbert} — substituição de URLs e menções, remoção de emojis, remoção do símbolo de \textit{hashtag}, normalização de espaços e normalização de caixa com preservação de entidades financeiras;
	
	\item \textbf{Carregamento do modelo}: o tokenizador e o modelo são instanciados via \texttt{AutoToke nizer} e \texttt{BertForSequenceClassification}, encapsulados em um objeto \texttt{pipeline} da biblioteca \texttt{transformers} com \texttt{task="text-classification"};
	
	\item \textbf{Inferência}: cada publicação é classificada pelo método \texttt{predict\_text()}, com \texttt{batch\_size=32}. O modelo retorna o \textit{label} predito e o escore de confiança (\textit{score}) associado;
	
	\item \textbf{Mapeamento de \textit{labels}}: os rótulos em inglês retornados pelo modelo são normalizados para o português: \texttt{POSITIVE} → \texttt{positivo}, 	\texttt{NEGATIVE} → \texttt{negativo}, \texttt{NEUTRAL} → \texttt{neutro};
	
	\item \textbf{Persistência e visualização}: os resultados são armazenados em \texttt{tweets\_classifi cation} com \texttt{classificator = "FinBERT-PT-BR"} e o respectivo 	\texttt{score}. O painel Streamlit de processamento (\texttt{app/dashboard/pages/processing.py}), invocado via \texttt{make dashboard}, apresenta a comparação entre o sentimento predito pelo modelo e o sentimento anotado pelo humano tweet a tweet, além de métricas de concordância.
\end{enumerate}

\subsection{Abordagem Baseada em Modelo de Linguagem Pré-treinado com \textit{Fine-tuning}: BERTimbau}
\label{subsec:bertimbau}

\subsubsection{Fundamentação da Escolha}
\label{subsubsec:justificativa_bertimbau}

O \textbf{BERTimbau} \cite{souza2020} foi selecionado como modelo base para a abordagem de \textit{fine-tuning} local por três razões complementares:

\begin{enumerate}
	\item \textbf{Estado da arte para o português brasileiro genérico}: o BERTimbau é o modelo BERT pré-treinado de referência para a língua portuguesa, treinado sobre um corpus de 2,68 bilhões de palavras extraídas de fontes diversificadas --- Wikipedia em português, Common Crawl e BrWaC \cite{souza2020}. Sua adoção como base de \textit{fine-tuning} é amplamente consolidada na literatura de PLN em português, inclusive como ponto de partida do próprio FinBERT-PT-BR \cite{santos2023finbert};
	
	\item \textbf{Endereçamento direto do desalinhamento de domínio}: o FinBERT-PT-BR foi pré-ajustado sobre notícias e relatórios financeiros, gênero textual estruturalmente distinto das publicações do X/Twitter. O \textit{fine-tuning} do BERTimbau diretamente sobre o corpus de \textit{tweets} coletados --- incluindo suas características estilísticas de informalidade, densidade de abreviações e uso de emojis e \textit{hashtags} --- aborda essa limitação de forma explícita, produzindo um modelo alinhado ao gênero e ao subdomínio específico avaliado;
	
	\item \textbf{Adequação ao tamanho do corpus}: a versão \texttt{base} do BERTimbau (110 milhões de parâmetros) foi preferida à versão \texttt{large} (335 milhões de parâmetros) pela maior adequação ao tamanho do corpus rotulado disponível. Modelos de maior capacidade requerem volumes de dados proporcionalmente maiores para generalização estável, e o uso da versão \texttt{large} sobre um corpus de algumas centenas ou poucos milhares de exemplos aumentaria o risco de \textit{overfitting} sem ganho garantido de desempenho \cite{devlin2019}.
\end{enumerate}

\subsubsection{Arquitetura e Protocolo de \textit{Fine-tuning}}
\label{subsubsec:arquitetura_bertimbau}

O modelo base utilizado é o \texttt{neuralmind/bert-base-portuguese-cased}, disponível publicamente no repositório \textit{Hugging Face}. Sobre esse modelo, realiza-se o \textit{fine-tuning} para a tarefa de classificação de sequência em três classes --- \textbf{Positivo}, \textbf{Neutro} e \textbf{Negativo} ---, por meio da adição de uma camada de classificação linear sobre a representação do \textit{token} especial \texttt{[CLS]}, que codifica o sentido global da sequência de entrada \cite{devlin2019}.

O protocolo de \textit{fine-tuning} é implementado no script \texttt{app/core/processing/bert /bert\_timbau\_fine\_tuner.py} e compreende as seguintes decisões de configuração:

\begin{itemize}

	\item \textbf{Particionamento}: o corpus rotulado é particionado em etapa dedicada, anterior ao \textit{fine-tuning}, conforme o protocolo descrito na \autoref{subsec:particionamento}. O particionamento é gerenciado pela classe \texttt{DatasetSplitRepository} (\texttt{app/shared/db/dataset\_split.py}), que persiste o resultado na tabela \texttt{dataset\_split} em dois conjuntos mutuamente exclusivos: (i)~\textbf{hold-out de teste} ($\approx$~15\% dos registros), extraído por amostragem estratificada com \texttt{random\_state=42} e reservado exclusivamente para a avaliação comparativa final descrita na \autoref{subsec:metricas_avaliacao}, marcado com \texttt{split=`test'} e \texttt{fold=NULL}; e (ii)~\textbf{conjunto de treinamento} ($\approx$~85\%), subdividido em $K=4$ folds por validação cruzada estratificada (\texttt{StratifiedKFold}, $K=4$, \texttt{shuffle=True}, \texttt{random\_state=42}), marcado com \texttt{split=`train'} e \texttt{fold}~$\in \{1, 2, 3, 4\}$;
	
	\item \textbf{Validação cruzada e seleção do modelo}: o \textit{fine-tuning} é executado $K=4$ vezes, uma por fold. Em cada iteração~$k$, os registros com \texttt{fold}~$\neq k$ ($\approx$~64\% do total rotulado) formam o conjunto de treinamento efetivo, e os registros com \texttt{fold}~$= k$ ($\approx$~21\%) formam o conjunto de validação. Ao término das $K$ iterações, o modelo correspondente ao fold com maior F1-Score macro na validação é selecionado como modelo definitivo e persistido em \texttt{models/bert-timbau-sentiment/}, acompanhado do relatório de classificação e do resumo do K-Fold em \texttt{models/bert-timbau-sentiment/eval/};
	
	\item \textbf{Ponderação de classes}: dado o desbalanceamento esperado em favor da classe neutra, os pesos por classe são calculados via \texttt{compute\_class\_weight(class\_weight=`balanced')} da biblioteca \texttt{scikit-learn} e injetados na função de perda (\texttt{CrossEntropyLoss}) por meio de uma subclasse \texttt{WeightedTrainer}, garantindo que o modelo não colapse para predizer sistematicamente a classe majoritária;
	
	\item \textbf{Hiperparâmetros}: taxa de aprendizado $5 \times 10^{-5}$, tamanho de lote (\textit{batch size}) 8, comprimento máximo de sequência 512 \textit{tokens}, razão de \textit{warmup} 0,1 e decaimento de pesos (\textit{weight decay}) $10^{-2}$. Esses valores foram determinados por busca empírica sobre o conjunto de validação cruzada e são consistentes com as recomendações dos autores do BERT para \textit{fine-tuning} em conjuntos de dados de tamanho moderado \cite{devlin2019};
	
	\item \textbf{\textit{Early stopping}}: o treinamento de cada fold é limitado a doze épocas, com critério de parada antecipada (\textit{patience}~=~2) monitorando o F1-Score macro no conjunto de validação do fold corrente. O \textit{checkpoint} com melhor F1-Score macro é restaurado automaticamente ao término de cada iteração, prevenindo o \textit{overfitting} à distribuição de treinamento;
	
	\item \textbf{Reprodutibilidade}: a semente aleatória \texttt{random\_state=42} é fixada em todas as operações estocásticas do treinamento (\texttt{seed} do \texttt{TrainingArguments}), garantindo que execuções sucessivas sobre o mesmo corpus produzam resultados idênticos.
\end{itemize}

O modelo ajustado, o tokenizador e o mapeamento de rótulos (\texttt{id2label.json}) são persistidos em \texttt{models/bert-timbau-sentiment/}, e o relatório de classificação e a matriz de confusão no conjunto de validação são salvos em \texttt{models/bert-timbau-sentiment/eval/}, para referência na análise apresentada no \autoref{Cap:Ava_Exp}.

\subsubsection{Implementação e Protocolo de Inferência}
\label{subsubsec:protocolo_inferencia_bertimbau}

A inferência é implementada pela classe \texttt{BERTimbauAnalyzer}, definida em \texttt{app/core/ processing/bert/bert\_timbau.py}. Assim como \texttt{FinBertPTBRAnalyzer}, a classe herda de \texttt{BaseSentimentAnalyzer} (\texttt{app/core/processing/base\_analyzer.py}), garantindo que ambos os analisadores compartilhem o mesmo contrato de interface e possam ser alternados pelo painel de processamento sem modificações adicionais. O protocolo de inferência é análogo ao descrito na \autoref{subsubsec:protocolo_inferencia}, com as seguintes particularidades:

\begin{enumerate}
	\item \textbf{Pré-requisito de treinamento}: a classe verifica na inicialização a existência do diretório \texttt{models/bert-timbau-sentiment/}; caso ausente, emite um erro orientando o usuário a executar \texttt{python -m pipeline.05\_processing.bert\_timbau\_fine\_tuner} (atalho: \texttt{make finetune}) antes da inferência;
	
	\item \textbf{Carregamento do modelo}: tokenizador e modelo são carregados diretamente do diretório local via \texttt{AutoTokenizer} e \texttt{AutoModelForSequenceClassification}, encapsulados em objeto \texttt{pipeline} da biblioteca \texttt{transformers} com \texttt{task="text-classifi cation"};
	
	\item \textbf{Mapeamento de rótulos}: o modelo fine-tuned emite rótulos em português (\texttt{positivo}, \texttt{negativo}, \texttt{neutro}) diretamente via \texttt{id2label}, dispensando a etapa de normalização de inglês para português aplicada ao FinBERT-PT-BR;
	
	\item \textbf{Persistência}: os resultados são armazenados em \texttt{tweets\_classification} com \texttt{classificator = "BERTimbau"} e o respectivo escore de confiança, seguindo o mesmo esquema de persistência descrito na \autoref{subsubsec:protocolo_inferencia} para o FinBERT-PT-BR.
\end{enumerate}


\subsection{Métricas de Avaliação das Abordagens}
\label{subsec:metricas_avaliacao}

A comparação entre as quatro abordagens --- OpLexicon, SentiLex-PT, FinBERT-PT-BR e BERTimbau ajustado --- será realizada exclusivamente sobre o conjunto de teste \textit{hold-out} definido na \autoref{subsec:particionamento}, que constitui o \textit{gold standard} efetivo da avaliação. Essa restrição é necessária para garantir comparabilidade: como o BERTimbau ajustado viu parte do corpus rotulado durante o treinamento, avaliá-lo sobre o conjunto completo produziria estimativas otimistas de desempenho incomparáveis às das demais abordagens. A implementação das métricas está organizada no sexto estágio do \textit{pipeline} (\texttt{app/core/evaluation/metrics.py}), invocado via \texttt{make evaluate}. As métricas calculadas são:

\begin{itemize}
	\item \textbf{Acurácia}: proporção de classificações corretas sobre o total de instâncias do conjunto de teste;
	
	\item \textbf{Precisão, Revocação e F1-Score} por classe: métricas calculadas individualmente para as classes Positivo, Negativo e Neutro, permitindo identificar vieses de classificação em classes específicas;
	
	\item \textbf{F1-Score macro}: média aritmética não ponderada dos F1-Scores de cada classe, tratando as classes como igualmente importantes independentemente de seu tamanho --- métrica de referência para comparação entre abordagens em corpora desbalanceados;
	
	\item \textbf{Matriz de confusão}: representação tabular da distribuição de acertos e erros por classe, útil para identificar padrões sistemáticos de classificação incorreta entre as três abordagens.
\end{itemize}

A limitação inerente à anotação por avaliador único --- ausência de cálculo de concordância interanotadores (\textit{inter-annotator agreement}) --- é reconhecida como uma restrição metodológica do presente trabalho.